\documentclass[11pt]{article}
\usepackage[utf8]{inputenc}
\usepackage{amsmath, amssymb, amsthm, mathtools, bm, geometry, setspace}
\geometry{margin=1in}
\onehalfspacing

% --- Independent Reindexing for Each Category ---
\newtheorem{assumption}{Assumption}
\newtheorem{proposition}{Proposition}
\newtheorem{theorem}{Theorem}
\newtheorem{lemma}{Lemma}
\newtheorem{definition}{Definition}
\newtheorem{corollary}{Corollary}

\title{Auditing Opaque Platforms via Incentivized Switching: A Partial Identification Approach}
\author{Qifan Han}
\date{March 2026}

\begin{document}
\maketitle

\begin{abstract}
We study the problem of identifying the causal effect of algorithmic treatments (e.g., digital ads) when the platform's assignment rule is opaque. Standard experimental designs fail to identify the Average Treatment Effect (ATE) because the platform's latent targeting induces one-sided noncompliance. While recent literature proposes active multi-cell designs to recover these effects, third-party auditors often lack the platform cooperation required to implement them. We propose an "Incentivized Switching" design where the auditor randomly assigns financial incentives to users to bypass the platform's logic. Building on the robust partial identification literature, we characterize the sharp identification region for the ATE without parametric assumptions on the unobserved propensity score. We introduce the \textit{Selection Ratio} to quantify the extent of algorithmic cherry-picking.
\end{abstract}

\section{Introduction}
In digital ecosystems, platforms (e.g., Google, Meta, OpenAI) deploy proprietary algorithms to assign treatments such as advertisements or content recommendations. For an external econometrician or auditor, measuring the causal impact of these interventions poses a fundamental challenge: the platform's targeting criteria are unobserved, and the assignment is correlated with potential outcomes.

As highlighted by Waisman and Gordon (2025), standard A/B tests in these environments suffer from one-sided noncompliance. If an advertiser restricts a campaign to a randomized eligibility group, the platform still retains ultimate control over which eligible users actually receive the ad. Consequently, standard estimators recover only the Average Treatment Effect on the Treated (ATT) for the platform's latently selected subpopulation. To solve scale-up and intensive-margin decisions (i.e., identifying the Average Treatment Effect, ATE), Waisman and Gordon (2025) demonstrate the necessity of multi-cell experiments that actively manipulate the platform's bid landscape. 

However, third-party auditors, researchers, or smaller advertisers cannot force platforms to execute multi-cell designs. They face a rigid, black-box assignment rule. How can an auditor measure the ATE when the platform refuses to cooperate?

We propose an **Incentivized Switching** experimental design. Consider a platform with a free tier (where ads are shown based on the black-box algorithm) and a paid tier (where ads are strictly suppressed). The auditor recruits a sample of users and randomly assigns financial incentives to upgrade to the paid tier. This guarantees a clean counterfactual state. Yet, because the auditor cannot observe which users in the free tier were latently targeted by the platform, point identification of the ATE remains impossible.

To resolve this, we map this institutional friction into the robust partial identification framework developed by Manski (1990, 1997) and extended in recent literature on unobserved heterogeneity (e.g., Kaido, Molinari, and Stoye 2019; Barseghyan et al. 2021). Rather than imposing ad-hoc parametric assumptions on the unobserved platform propensity score to force point identification, we maintain a flexible latent-index model and characterize the sharp identification region for the ATE. By imposing a theoretically motivated shape restriction---Monotone Selection on Gains---we demonstrate how auditors can systematically bound the platform's true causal effect and quantify its algorithmic targeting efficiency.

\section{The Model and Experimental Setup}
We consider a population of users indexed by $i$. The platform administers a binary treatment. Let $D_i \in \{0,1\}$ denote the platform's latent intent to treat user $i$. 

The auditor implements an experiment by recruiting users and randomly assigning an incentive $Z_i \in \{0,1\}$. 
\begin{itemize}
    \item $Z_i = 1$: The user is incentivized to switch to a "Premium/Paid" tier where the platform's treatment is technically suppressed.
    \item $Z_i = 0$: The user remains in the "Free" tier subject to the platform's natural targeting $D_i$.
\end{itemize}

Let $W_i \in \{0,1\}$ denote the realized treatment receipt. The institutional constraints dictate that $W_i = D_i(1-Z_i)$. 

Each user has potential outcomes $Y_i(w)$ for $w \in \{0,1\}$. The observed outcome is:
\begin{equation}
    Y_i = W_i Y_i(1) + (1-W_i) Y_i(0)
\end{equation}
The individual treatment effect is $\tau_i = Y_i(1) - Y_i(0)$. The auditor observes the vector $(Y_i, Z_i, X_i)$, where $X_i$ are baseline covariates. Crucially, in many auditing scenarios, the auditor does not reliably observe $W_i$ or $D_i$.

\begin{assumption}[Randomized Auditor Intervention]
\label{ass:randomization}
Conditional on observables $X_i$, the auditor's incentive assignment $Z_i$ is independent of the platform's latent intent and the user's potential outcomes:
\[ Z_i \perp (D_i, Y_i(0), Y_i(1)) \mid X_i \]
\end{assumption}

Assumption \ref{ass:randomization} ensures that the recruitment and incentivization process does not suffer from endogeneity.

\section{Identification of the Diluted Effect}
Before establishing bounds on the ATE, we first formally characterize what the incentivized experiment point-identifies. Let $p(x) = P(D_i=1 \mid X_i=x)$ denote the unobserved platform propensity score.

\begin{proposition}[Identification of the Diluted ATT]
\label{prop:diluted}
Under Assumption \ref{ass:randomization}, the difference in conditional expectations between the control and treated arms of the auditor's experiment point-identifies the product of the unobserved propensity score and the ATT:
\begin{equation}
    E[Y_i \mid Z_i=0, X_i=x] - E[Y_i \mid Z_i=1, X_i=x] = p(x) \cdot ATT(x)
\end{equation}
where $ATT(x) = E[\tau_i \mid D_i=1, X_i=x]$.
\end{proposition}

Proposition \ref{prop:diluted} highlights the core identification failure: because the auditor cannot observe $p(x)$, they cannot disentangle the treatment effect ($ATT$) from the platform's targeting coverage ($p$).

\section{Partial Identification of the ATE}
Following the econometric tradition of Barseghyan et al. (2021), we avoid unverifiable exclusion restrictions. Instead, we formalize the platform's assignment mechanism and place shape restrictions on the heterogeneity of treatment effects to construct a sharp identification region for the $ATE(x) = E[\tau_i \mid X_i=x]$.

\begin{assumption}[Latent Index Assignment]
\label{ass:latent}
There exists a scalar latent variable $U_i$ such that $D_i = \mathbf{1}\{p(X_i) \ge U_i\}$, with $U_i \mid X_i \sim \mathrm{Unif}[0,1]$.
\end{assumption}

We define the Marginal Treatment Effect as $MTE(x, u) = E[\tau_i \mid X_i=x, U_i=u]$.

\begin{assumption}[Monotone Selection on Gains]
\label{ass:monotone}
For any $x$, the function $u \mapsto MTE(x, u)$ is weakly decreasing on $[0,1]$.
\end{assumption}

Assumption \ref{ass:monotone} reflects the economic incentive of the platform: it prioritizes treatment for users who exhibit the highest marginal response.

\begin{assumption}[Non-Negative Treatment Effects]
\label{ass:nonneg}
For any $(x, u)$, $MTE(x, u) \ge 0$.
\end{assumption}

Let $\Delta(x) \equiv E[Y_i \mid Z_i=0, X_i=x] - E[Y_i \mid Z_i=1, X_i=x]$ denote the observable experimental difference. Let $\Theta_I(x)$ denote the sharp identification region for $ATE(x)$.

\begin{theorem}[Sharp Identification Region for ATE]
\label{thm:sharpbounds}
Suppose Assumptions \ref{ass:randomization}--\ref{ass:nonneg} hold. Further suppose the auditor has access to a lower bound on the platform's propensity score, $\underline{p}(x) \le p(x)$, obtained from external industry knowledge or bounds on sample observables. The sharp identification region for $ATE(x)$ is:
\begin{equation}
    \Theta_I(x) = \left[ \Delta(x), \,\, \frac{\Delta(x)}{\underline{p}(x)} \right]
\end{equation}
\end{theorem}

\section{Auditing the Black Box: The Selection Ratio}
The theoretical bounds in Theorem \ref{thm:sharpbounds} yield a practical auditing metric. We define the \textbf{Selection Ratio} as the degree to which the platform's targeted ROI exceeds the population ROI:
\begin{equation}
    R(x) = \frac{ATT(x)}{ATE(x)}
\end{equation}
If the platform assigns treatments randomly, $R(x)=1$. If the algorithm is highly efficient at "cherry-picking," $R(x)$ is strictly greater than 1.

\begin{corollary}[Bounds on Platform Efficiency]
\label{cor:ratio}
Under the conditions of Theorem \ref{thm:sharpbounds}, the Selection Ratio is partially identified such that:
\begin{equation}
    R(x) \in \left[ 1, \,\, \frac{1}{\underline{p}(x)} \right]
\end{equation}
\end{corollary}

\section{Proofs}

\subsection*{Proof of Proposition \ref{prop:diluted}}
Based on the experimental setup, if $Z_i=1$, $W_i=0$ and thus $Y_i = Y_i(0)$. If $Z_i=0$, $W_i = D_i$ and thus $Y_i = D_i Y_i(1) + (1-D_i)Y_i(0)$. Taking expectations conditional on $X_i=x$:
\begin{align*}
    E[Y_i \mid Z_i=1, x] &= E[Y_i(0) \mid x] \\
    E[Y_i \mid Z_i=0, x] &= E[D_i Y_i(1) + (1-D_i)Y_i(0) \mid x] = E[Y_i(0) \mid x] + E[D_i(Y_i(1)-Y_i(0)) \mid x]
\end{align*}
where we utilize Assumption \ref{ass:randomization} to ensure $E[Y_i(0) \mid Z_i=1, x] = E[Y_i(0) \mid Z_i=0, x]$. Subtracting the expectations yields the observable difference $\Delta(x) = E[D_i \tau_i \mid x]$. By the definition of conditional probability:
\[ E[D_i \tau_i \mid x] = P(D_i=1 \mid x) E[\tau_i \mid D_i=1, x] = p(x) \cdot ATT(x) \]
\qed

\subsection*{Proof of Theorem \ref{thm:sharpbounds}}
By Assumption \ref{ass:latent}, we rewrite the ATE and ATT using the MTE:
\[ ATE(x) = \int_0^1 MTE(x,u) du = p(x)ATT(x) + (1-p(x))ATU(x) \]
where $ATT(x) = \frac{1}{p(x)}\int_0^{p(x)} MTE(x,u) du$ and $ATU(x) = \frac{1}{1-p(x)}\int_{p(x)}^1 MTE(x,u) du$.
\newline
\textbf{Lower Bound:} By Assumption \ref{ass:nonneg}, $ATU(x) \ge 0$. Therefore, $ATE(x) \ge p(x)ATT(x)$. From Proposition \ref{prop:diluted}, $p(x)ATT(x) = \Delta(x)$. Thus, $ATE(x) \ge \Delta(x)$.
\newline
\textbf{Upper Bound:} By Assumption \ref{ass:monotone}, the average of the weakly decreasing MTE over $[0, p(x)]$ must be weakly greater than the average over $[p(x), 1]$. Thus, $ATU(x) \le ATT(x)$, implying $ATE(x) \le ATT(x)$. Since $ATT(x) = \Delta(x)/p(x)$ and $p(x) \ge \underline{p}(x)$, it follows that $ATE(x) \le \Delta(x)/\underline{p}(x)$.
\newline
\textbf{Sharpness:} To prove sharpness, we must show that for any $\theta \in [\Delta(x), \Delta(x)/\underline{p}(x)]$, there exists a data generating process satisfying all assumptions that yields $ATE(x) = \theta$. 
For the lower bound, let $MTE(x,u) = \Delta(x)/p(x)$ for $u \le p(x)$ and $0$ otherwise. This satisfies monotonicity and non-negativity, and yields $ATE(x) = \Delta(x)$.
For the upper bound, let $MTE(x,u) = \Delta(x)/p(x)$ for all $u \in [0,1]$. This is a flat, non-negative function satisfying monotonicity, and yields $ATE(x) = \Delta(x)/p(x)$. Varying $p(x)$ down to $\underline{p}(x)$ attains the supremum. Since convex combinations of these MTE schedules are also valid, the entire interval is attainable. \qed

\subsection*{Proof of Corollary \ref{cor:ratio}}
By definition, $R(x) = ATT(x)/ATE(x)$. 
Using the bounds from Theorem \ref{thm:sharpbounds}:
The minimum ratio occurs when $ATE(x)$ is maximized (i.e., $ATE(x) = ATT(x)$), yielding $R(x) = 1$.
The maximum ratio occurs when $ATE(x)$ is minimized (i.e., $ATE(x) = \Delta(x)$). Since $ATT(x) = \Delta(x)/p(x)$, the maximum ratio is $1/p(x)$. Since $p(x) \ge \underline{p}(x)$, the upper bound on the ratio is $1/\underline{p}(x)$. \qed

\end{document}